\section{The exact law}\label{sec:exact}

The results of this section are classical~\cite{renshaw1981,dekkingkong2011}.
We prove Theorem~\ref{thm:pmf} because the later sections use its count of
words by switches. By \eqref{eq:wordprob}, a word with $j$ switches has
probability $\frac12p^{N-1-j}q^j$.

\begin{theorem}[bin probabilities]\label{thm:pmf}
Let $N\ge1$. Then $P_N(0)=P_N(N)=\tfrac12p^{N-1}$, and for $1\le k\le N-1$
\[
P_N(k)=\frac12\sum_{j=1}^{N-1}W(N,k,j)\,p^{N-1-j}q^{\,j},
\]
where, writing $a=k$ and $b=N-k$,
\[
W(N,k,2i-1)=2\binom{a-1}{i-1}\binom{b-1}{i-1},\qquad
W(N,k,2i)=\binom{a-1}{i}\binom{b-1}{i-1}+\binom{a-1}{i-1}\binom{b-1}{i}
\]
is the number of words with $a$ letters $\sR$, $b$ letters $\sL$ and exactly
$j$ switches.
\end{theorem}

\begin{proof}
By \eqref{eq:wordprob} it suffices to count the words in bin $k$ with $j$
switches. We count them by their run lengths. A word with $j$ switches
consists of $j+1$ maximal runs, alternately of letters $\sR$ and of letters
$\sL$. If $k\in\{0,N\}$ there is one word and it has no switch. Let $a,b\ge1$.
If $j+1=2i$ there are $i$ runs of each letter. The run lengths form a
composition of $a$ into $i$ parts and a composition of $b$ into $i$ parts, and
there are $\binom{a-1}{i-1}\binom{b-1}{i-1}$ such pairs. The two choices of the
first letter give the factor $2$. If
$j+1=2i+1$ the word starts and ends with the same letter. Starting with $\sR$
gives $i+1$ runs of $\sR$ and $i$ runs of $\sL$, so
$\binom{a-1}{i}\binom{b-1}{i-1}$ words, and starting with $\sL$ gives the
other summand.
\end{proof}

For instance $P_2=(\tfrac p2,\,q,\,\tfrac p2)$, $P_3(1)=\tfrac12(1-p^2)$ and
$P_4(2)=q(1-p+p^2)$. Each interior bin needs at least one switch, while the
endpoints need none. So the endpoints win when $p$ is close to $1$, and
Section~\ref{sec:middle} asks how close.

\begin{proposition}[recurrence and generating function]\label{prop:gf}
Let $P_N^{\sR}(k)$ and $P_N^{\sL}(k)$ be the probabilities of being in bin $k$
after $N$ steps with last letter $\sR$, respectively $\sL$, and set them to $0$
outside $0\le k\le N$. Then $P_1^{\sR}(1)=P_1^{\sL}(0)=\tfrac12$,
$P_1^{\sR}(0)=P_1^{\sL}(1)=0$, and for $N\ge2$
\begin{equation}\label{eq:dp}
P_N^{\sR}(k)=p\,P_{N-1}^{\sR}(k-1)+q\,P_{N-1}^{\sL}(k-1),\qquad
P_N^{\sL}(k)=p\,P_{N-1}^{\sL}(k)+q\,P_{N-1}^{\sR}(k).
\end{equation}
Consequently
\begin{equation}\label{eq:gf}
\sum_{N\ge0}\sum_{k=0}^{N}P_N(k)\,x^kt^N=\frac{1-(p-\tfrac12)(1+x)\,t}{1-p(1+x)\,t+(2p-1)\,x\,t^2},
\end{equation}
and for all $N\ge2$ and all $k$, with $P_M(k)=0$ outside $0\le k\le M$,
\begin{equation}\label{eq:rec}
P_N(k)=p\bigl(P_{N-1}(k)+P_{N-1}(k-1)\bigr)-(2p-1)\,P_{N-2}(k-1).
\end{equation}
\end{proposition}

\begin{proof}
The recursion \eqref{eq:dp} follows by conditioning on $w_{N-1}$, since $w_N$
equals $w_{N-1}$ with probability $p$.
Let $H^{\sR}=\sum_{N,k}P_N^{\sR}(k)x^kt^N$ and define $H^{\sL}$ alike. By
\eqref{eq:dp}, $H^{\sR}=xt(\tfrac12+pH^{\sR}+qH^{\sL})$ and
$H^{\sL}=t(\tfrac12+pH^{\sL}+qH^{\sR})$. Solving these two linear equations
and adding $1$ for $N=0$ gives \eqref{eq:gf}, which is the generating function
of Renshaw and Henderson~\cite{renshaw1981} in our notation. Multiplying
\eqref{eq:gf} by its denominator and comparing the coefficients of $x^kt^N$
with $N\ge2$ gives \eqref{eq:rec}.
\end{proof}

\begin{remark}[numerator tables]\label{rem:tables}
For $p=\alpha/(\alpha+\beta)$ with integers $\alpha,\beta\ge0$, not both $0$,
the numerators of the tables of the problem are
$n_N(k)=2(\alpha+\beta)^{N-1}P_N(k)$ for $N\ge1$, not reduced. By
\eqref{eq:wordprob}, $n_N(k)$ is the sum of $\alpha^{\rep(w)}\beta^{\chg(w)}$
over the words with $k$ letters $\sR$ and $N-k$ letters $\sL$, which shows that
it is an integer.
\end{remark}

\subsection{Variance and the normal limit}\label{sec:clt}

Note that $K_N$ counts the visits of the chain $w_1,\dots,w_N$ to the letter
$\sR$. This chain leaves each letter with probability $q$ and starts from its
stationary law $(\frac12,\frac12)$. So $K_N$ has the Markov binomial law
$\mathrm{Bin}(N,q,q,(\frac12,\frac12))$ of Dekking and
Kong~\cite{dekkingkong2011}, whose second eigenvalue is $1-2q=2p-1$.

\begin{proposition}[mean and variance]\label{prop:var}
Let $0<p<1$ and $\rho=2p-1$. Then $\E K_N=N/2$ and
\[
\Var K_N=\frac14\Bigl(N\,\frac{1+\rho}{1-\rho}-\frac{2\rho\,(1-\rho^N)}{(1-\rho)^2}\Bigr)
=\frac{Np}{4(1-p)}-\frac{(2p-1)\bigl(1-(2p-1)^N\bigr)}{8(1-p)^2}.
\]
\end{proposition}

\begin{proof}
This is \cite[Eq.~(2.3) and Proposition~2.1]{dekkingkong2011} with $a=b=q$
and the stationary start, for which both excentricities vanish and both
stationary weights are $\frac12$. The second form follows from
$(1+\rho)/(1-\rho)=p/q$ and $1-\rho=2q$.
\end{proof}

\begin{theorem}[normal limit]\label{thm:clt}
Let $0<p<1$ be fixed. Then $(K_N-N/2)/\sqrt N$ converges in distribution, as
$N\to\infty$, to the normal law with mean $0$ and variance $p/(4(1-p))$.
\end{theorem}

So persistence keeps the bell shape but widens it by the factor $\sqrt{p/q}$
compared with the ordinary board.

\begin{proof}
Let $v_N(\theta)$ be the vector with entries
$\E[e^{i\theta X_N};\sigma_N=\pm1]$. By \eqref{eq:dp}, $v_{N+1}=Mv_N$ with
\[
M=\begin{pmatrix}pe^{i\theta}&qe^{i\theta}\\qe^{-i\theta}&pe^{-i\theta}\end{pmatrix},\qquad
\operatorname{tr}M=2p\cos\theta,\quad\det M=2p-1 .
\]
At $\theta=0$ the eigenvalues are $1$ and $2p-1$, which are distinct. So for
small $|\theta|$ the matrix $M$ has the simple eigenvalue
$\lambda_+(\theta)=p\cos\theta+\sqrt{q^2-p^2\sin^2\theta}
=1-\frac p{2q}\theta^2+O(\theta^4)$ and the other eigenvalue stays near $2p-1$,
with spectral projections $\Pi_\pm(\theta)$ that depend continuously on
$\theta$, because the two eigenvalues are simple. Write $\mathbf 1=(1,1)$. Then
$\E e^{i\theta X_N}=\mathbf 1\cdot M^{N-1}v_1(\theta)
=\lambda_+^{N-1}\,\mathbf 1\cdot\Pi_+v_1+\lambda_-^{N-1}\,\mathbf 1\cdot\Pi_-v_1$.
The vector $v_1(0)=(\frac12,\frac12)$ is an eigenvector of $M$ for the
eigenvalue $1$ at $\theta=0$, so $\mathbf 1\cdot\Pi_+v_1\to1$ as $\theta\to0$,
while $|\lambda_-|\to|2p-1|<1$. Put $\theta=s/\sqrt N$ with $s$ fixed. Then
the second term tends to $0$, and the first tends to $\exp(-\frac p{2q}s^2)$
since $(N-1)\log\lambda_+(s/\sqrt N)\to-\frac p{2q}s^2$. By
L\'evy's continuity theorem
\cite[Theorem~26.3]{billingsley1995}, $X_N/\sqrt N$ converges to the normal
law with variance $p/q$. Finally $K_N-N/2=X_N/2$.
\end{proof}
